\section{本章小结}
\label{sec:bo-summary}

本章的出发点是预算极少的黑箱问题，$T$次评价要在有希望的区域和未知的区域之间分配。评价点负责收集信息，推荐解负责交付，选点规则只能使用已完成的观测。目标本身并不随机，概率描述的是我们对它的认识。

概率代理把已完成的观测变成对所有候选位置的预测。高斯过程用核表达“相近的位置函数值相近”，一个读数通过后验协方差影响相关的位置，后验方差只取决于评价放在哪里。超参数不确定时，直接代入估计值可能严重低估不确定性，低方差只说明在模型内不确定性小。

采集函数把预测变成分数：PI只看过线的概率，EI把幅度也计入，LCB直接按乐观值比较，Thompson采样按后验抽样选点，KG衡量一次观测对最终推荐的改善。四种规则的共同之处是都在利用与探索之间折中，不同之处在于如何折中。循环的外层是昂贵评价，内层是对采集函数的优化，而内层的求解质量也影响选点。

各种推广分别改变了循环的某一环。噪声改变了阈值的含义，需要区分潜在函数值和下一次读数，复测的价值要由信息更新的准则衡量。约束在采集分数中加入可行概率，$b_t$也改为可行读数中的最好值；要求每次试验都可行时，还需要额外的安全假设。并行使一批候选之间的相关性变得重要，未完成的试验要按预测分布平均，而不当作已有读数。成本则改变了预算的口径。

最后要分开三件事：模型给出的概率是否可信，优化的结果是否够好，计算的开销是否值得。特定的GP置信界方法有遗憾上界，但它依赖模型正确、搜索域有限和置信系数的具体取值；这些条件不满足时，算法仍可以使用，只是结论要靠实际的对照比较来支持。
